% !TeX spellcheck = en_US
\documentclass[12pt]{article}

\input{common/variables}
\input{common/dates}

\setcounter{secnumdepth}{4}

\usepackage{times,fullpage,xspace,fancyhdr,url,color}
\usepackage[pdftex]{graphicx}
\usepackage[pdftex,
            colorlinks=true,
            urlcolor=black,
            linkcolor=black,
            citecolor=black,
            bookmarksopen=false,
            bookmarksnumbered=true,
            pdfstartview=FitH]{hyperref}

\pdfcompresslevel=9

\hypersetup{
 pdftitle={\leaguename (\leaguenameabbr) Technical Challenges},
 pdfauthor={Technical Committee \leaguenameabbr},
}
\usepackage{microtype}
\usepackage[utf8]{inputenc}
\usepackage{amsmath}
\usepackage{xargs}
\usepackage[colorinlistoftodos,prependcaption,textsize=tiny]{todonotes}
\usepackage{siunitx}
\usepackage[capitalize,noabbrev]{cleveref}
\usepackage[official]{eurosym}
\usepackage[useregional]{datetime2}
\usepackage{subcaption}
\usepackage{enumitem}
\usepackage{xcolor}
\DTMlangsetup[en-GB]{ord=raise,monthyearsep={,\space}}

% modern and easy to use tables tblr
\usepackage{tabularray}

\newcommandx{\unsure}[2][1=]{\todo[linecolor=red,backgroundcolor=red!25,bordercolor=red,#1]{#2}}
\newcommandx{\change}[2][1=]{\todo[linecolor=blue,backgroundcolor=blue!25,bordercolor=blue,#1]{#2}}
\newcommandx{\info}[2][1=]{\todo[linecolor=green,backgroundcolor=green!25,bordercolor=green,#1]{#2}}
\newcommandx{\improvement}[2][1=]{\todo[linecolor=orange,backgroundcolor=orange!25,bordercolor=orange,#1]{#2}}

% comment 'disable' in to disable all the todo notes :)
\usepackage
[
%disable
]{todonotes}

\usepackage[theorems]{tcolorbox}
\newtcbtheorem[number within=section]{hintbox}{}%
{colback=red!10,colframe=red!45!black,fonttitle=\bfseries}{th}

\usepackage{tikz}
\usetikzlibrary{backgrounds,arrows,calc,math,positioning,shapes, decorations.markings,decorations.pathreplacing}

\usepackage{tikz-3dplot}

\usepackage{tocloft}
\cftsetindents{section}{0em}{2em}
\cftsetindents{subsection}{0em}{2em}
\renewcommand\cfttoctitlefont{\hfill\Large\bfseries}
\renewcommand\cftaftertoctitle{\hfill\mbox{}}

\sloppy
\newcommand{\ie}{\mbox{i.\,e.}\xspace}
\newcommand{\eg}{\mbox{e.\,g.}\xspace}
%\newcommand{\cf}{\mbox{cf.}\xspace}
\newcommand{\cf}{see\xspace}
% \newcommand{\comment}[1]{\marginpar{\pdfannot width 4in height .5in depth 8pt {/Subtype /Text /Contents (#1)}}}
\newcommand{\inparagraph}[1]{\paragraph{#1\hspace{-1em} }}


% some colors
\definecolor{orange}{rgb}{1,0.5,0}
\definecolor{red}{rgb}{1,0,0}
\definecolor{green}{rgb}{0,1,0}

\renewcommand{\theparagraph}{§~\thesubsubsection.\arabic{paragraph}}%


% proper names for the \Cref and \cref

%\crefname{lawtmp}{Law}{Laws}
%\Crefname{lawtmp}{Law}{Laws}

% simply use Paragraph numbers
\crefname{lawtmp}{§}{§}
\Crefname{lawtmp}{§}{§}
\Crefname{paragraph}{}{}

%%%%%%%%%%%%%%%%%%% TEXT FORMAT %%%%%%%%%%%%%%%%%%%%%%%

\setlength{\parindent}{0pt}
\setlength{\parskip}{12pt plus 6pt minus 3 pt}
\widowpenalty=10000
\clubpenalty=10000

% needed to align an image and text correctly side by side
%\newcommand{\imagebox}[1]{\raisebox{2ex}{\raisebox{-\height}{#1}}}


%%%%%%%%%%%%%%%%%%% TITLE %%%%%%%%%%%%%%%%%%%%%%%

\title{Technical Challenges\\\leaguename (\leaguenameabbr)}

\author{\leaguename Technical Committee\\\url{https://hsl.robocup.org/}}
\date{(DRAFT \RCYear Working Rules Document, as of \today)}

\renewcommand{\headrulewidth}{0.4pt}
\renewcommand{\footrulewidth}{0.4pt}

%%%%%%%%%%%%%%%%%%% PAGE HEADER FORMAT %%%%%%%%%%%%%%%%%%%%%%%

\pagestyle{fancy}
% make sure that a separate page can be identified 
\lhead{Technical Challenges}
\chead{\leaguename}
\rhead{\today}
\lfoot{}
%\cfoot{}
%\cfoot{\thepage}
\rfoot{}


\begin{document}

\pagenumbering{Roman}

\maketitle
\thispagestyle{empty}

\begin{center}
  Questions or comments on this document should be submitted via the following channels:
  \begin{itemize}
      \item \url{https://github.com/RoboCup-HumanoidSoccerLeague/HSL-Rules/issues};
      \item Discord server: \url{https://discord.gg/67upw6Nn}, channel \texttt{\#rule-book};
      \item Email: \url{hsl_tc@lists.robocup.org}
  \end{itemize}
\end{center}

\newpage
\setcounter{tocdepth}{1}
\tableofcontents

\thispagestyle{fancy}

\clearpage
\pagenumbering{arabic}
\setcounter{page}{1}

\newpage


\section{Introduction}
At RoboCup \RCYear, the Humanoid Soccer League will hold technical challenges, which are described in this document.
RoboCup \RCYear awards a trophy for winning challenges.
% and the option for pre-qualification if a team is not pre-qualified by other means.

Technical challenges are used in the HSL to develop technical capabilities which will be used in upcoming RoboCups in the main competition. 
The purpose is to give teams time to develop solutions and exchange ideas before they will be introduced into the main competition. 
Challenges are designed to move the league in a direction of further improvement of soccer skills and towards the overall goal of 2050. 
Each team is strongly encouraged to participate in these challenges to contribute to the league's advancement.

\subsection{Code Publication}
Every team participating in a challenge must publish the corresponding code used in that competition according to the HSL rule book, unless a specific challenge states otherwise.

%\subsection{Overall Ranking}
%The scores earned in each challenge may vary significantly in magnitude, so they must be scaled before computing the overall Technical Challenge ranking.
%Teams that do not participate in a challenge will receive 0 points for that challenge.
%For each challenge, the team with the highest raw score will receive 25 points, and the team with the lowest non-zero raw score will receive 5 points.
%A linear scaling function is then defined between these two points, and all other participating teams will receive a scaled score based on their raw performance according to this function.

%\include{challenges/fitness-tests}
%\newpage

% !TeX root = ../Challenges.tex
% !TeX spellcheck = en_US

\section{Open Research Challenge}
\label{ssec:orc}

The open research challenge is held to showcase innovations and research projects that support the development of the league beyond the core competition objectives. 

\subsection{Submission Requirements}
Teams must submit a brief description of their planned contribution by the day before the presentation at RoboCup 2026 to the Technical Committee (\url{hsl_tc@lists.robocup.org}).
Multiple contributions per team are allowed.
The description should include:
    \begin{itemize}
        \item \textbf{Objective}: The purpose and scope of the innovation.
        \item \textbf{Expected Impact}: How it benefits HSL or improves existing systems.
        \item \textbf{Implementation Feasibility}: Resources and time needed to accomplish/execute the contribution.
    \end{itemize}
    
\subsection{Presentation Format}
Participating teams have to prepare a short 10 minute oral presentation. The presentation must include a clear explanation of the contribution and its impacts. A live or video demonstration should be included if feasible.

Presented topics should focus on innovations beneficial to the league, but not necessarily related to winning matches. Examples include:

\begin{itemize}
\item Projects improving the safety and handling of robots
\item Projects that increase interoperability and collaboration
\item Projects that improve refereeing
\item Projects that enable game statistics and evaluation
\end{itemize}

Presentations are held publicly, encouraging members of all teams to join, even if they did not submit a contribution themselves.

\subsection{Evaluation Process}
The winner will be decided by a vote among the team leaders during the competition using the Condorcet method~\footnote{\url{https://en.wikipedia.org/wiki/Condorcet_method}}. Each participating team will vote for their top teams in order (excluding themselves). Teams are encouraged to evaluate the presentations based on the following criteria:

\begin{itemize}
\item \textbf{Relevance}: Alignment with HSL’s goals of growth and improvement.
\item \textbf{Novelty}: Originality and creativity of the idea.
\item \textbf{Feasibility}: Ease of application, practicality of implementation within league resources.
\item \textbf{Impact}: Potential benefit to the league or teams.
\item \textbf{Usability}: Ease of use, integration, and adaptability of the contribution.
\item \textbf{Live Demonstration}: Preference should be given to contributions that include a live demonstration.
\item \textbf{Documentation and Code Release}: Consideration of whether the code was made available beforehand for review and that it is well documented.
\end{itemize}

At a time decided by the designated referee, within one hour of the last demonstration if not otherwise specified, the captain of each team will submit the team’s rankings by filling out a form. Any points awarded by a team to itself will be disregarded.

\subsection{Full Code Release Requirement}
A \textbf{full code release} is mandatory for participating in this challenge, except for contributions that do not include any source code. The code must be publicly available by 2026-10-15
%no later than the \textbf{main code release deadline} from the HSL rule book 
and must be announced to the Technical Committee (\url{hsl_tc@lists.robocup.org}). Code must be well documented.
Ideally, the code should be released before the competition to allow interested teams and individuals to review and provide feedback!

\subsection{Publication}
The presentation slides, videos and all related material to the presented project must be made available to the Technical Committee on the day of the presentation. 
All presentations and their related code will be made publicly available on the HSL league website.
\newpage

\end{document}
